\section{Conclusiones}

El proyecto aborda un vacío documentado entre visión por computadora y educación
inclusiva: el reconocimiento de Braille escolar manuscrito en español, sin
conexión y en hardware de gama baja. La comparación de arquitecturas y
cuantización permitirá seleccionar una configuración con respaldo empírico y
evaluar su efecto en la retroalimentación docente.

\subsection{Limitaciones}
El corpus es reducido (40 páginas), la prueba externa reúne 5 páginas y la fase
pedagógica es cuasiexperimental, sin grupo control independiente. El alcance se
limita al Braille Grado 1 y a condiciones de iluminación rasante.
